\section{Related Work}
\label{sec:relwork}

% Goal: review current bx works, from lenses (get-based), to putlenses (put-based) and biflux.
%  adv: fully certified.  Emphasis on this point ! TOBE REVISED.
% Sketch of related work
%  lenses for tree first one
%  relational lenses
%  string lenses, matching lenses
%  point-free lenses
%  bidirectional graph transformation
%  putback-based approach
%  putlenses
%  biflux
%  pdl
%  finally, ours.

Foster et al.
propose the first bidirectional programming language lenses~\cite{Foster-lenses}
for synchronizing tree-structured data.
%
They design a series of lens combinators that
each one can be evaluated as either a forward function $get$ 
or a backward function $put$.
%, and proved these two functions together are well-behaved.
Large programs are composed through compositions of basic lenses.
The novelty of their work is by putting emphasis on types and totality
of lens transformations, and allow reasoning about totality 
and well-behavedness of lenses in a compositional way.
%
Following this idea, many lenses languages are proposed for different data format.
%
Bohannon et al.~\cite{Bohannon:06} propose relational lenses for manipulating relational data
using basic combinators $select$, $drop$ and $join$ and composition of them,
they also
design lenses~\cite{Bohannon:2008:BRL:1328438.1328487} for string data based on operations
on regular transducers and, and implement a type system for this lenses based on regular expressions.
Matching lenses~\cite{Barbosa:2010} starts considering  update strategies of the ``putback"
by generalizing the string lenses from key-based matching
to support a set of different alignment strategies.
Pacheco et al.~\cite{PachecoCunha:10} propose a point-free lenses
by defining suitable backward semantics for most of the standard point-free combinators.
Hidaka et al.~\cite{Hidaka:10} propose the first bidirectional programming language for dealing with graphs,
which designs the backward semantics for a graph algebra UnCAL.
Liu et al.~\cite{Liu:2007:BIX:1244381.1244386} gives a bidirectional interpretation of XQuery to make an XQuery expression can
be executed in two directions (forward and backward).
All the above existing approaches focus on
writing bidirectional programs from the $get$ direction,
provide a default backward transformation $put$
or limit control over the derived $put$ when defining the $get$.

Putlenses~\cite{Pacheco-putlenses} is the first putback-based lens language
which provides a set of putback-based lens combinators to let user write flexible ``put".
BiFluX~\cite{Pacheco-BiFluX} gives a new program by update paradigm
which inherently lets programmer to write putback function and
the system derives the forward $get$ for free.
PDL~\cite{HuPF14} is a treeless functional language for writing total $put$ functions,
such that existence of well-behavedness can be checked statically,
and a corresponding total $get$ can be derived automatically.
While none of them managed to fully prove the well-behavedness formally. % ?
In this paper, we design a minimalist putback-based core language
that is fully certified in \Agda\ and useful for building real putback-based languages such as BiFluX.


\section{Conclusion}
\label{sec:con}
% reification ?
Putback-based bidirectional transformation opens a new research area,
which aims to make BX to be deployed for really practical usages.
Thus not only allow specifying flexibility over the updates like existing putback-based languages,
but also need fully guarantee the correctness of derived $get$.
\BiGUL\ is a datatype-generic, and fully certified BX core language
which can be served as the base for other putback-based languages,
and it is indeed that BiFLuX can be translated into \BiGUL.
% extension
While there are still extensions can be made.
Currently alignment only handles the case that
elements with the same matching key in both source and view are unique
%, which indicates the update operations after alignment focus on single source/view element.
, but in fact there are situations that more than one elements have the same matching key value,
the result would be a list of elements for both source and view.
\BiGUL\ can be extended with iteration operation on list to handle this case.
%
% source dependency rearrangement already done.
% view dependency need new feature to handle it.
Another one is view dependency that means two view elements depend on each other.
Embedding these two view elements separately into source may lead to inconsistency, 
extending \BiGUL\ to support specifying the dependency between these two view elements is a good choice.
\todo[inline]{Towards a dependently typed version}

%For example the catalog of a book dependends on the strcuture and contents of the book.
%If the title of contents changes, the catalog shall be changed.
%This can be captured by expressing the depdenency between the catalog and the titles in the content.
%Only one of them (e.g. titles in the content) needs to be used to update the corresponding source element.
%From the $get$ perspective, catalog will be recovered through the view dependency.


